\section{深度学习的崛起}

在计算机视觉循序渐进发展的同时，一条平行的研究路线也在默默推进——这就是神经网络和后来的深度学习。

\subsection{神经网络的早期探索}

\subsubsection{感知机与 Minsky 的批评}

神经网络的历史可以追溯到 1958 年 Frank Rosenblatt 提出的\textbf{感知机}（Perceptron）。然而，1969 年 Marvin Minsky 和 Seymour Papert 证明感知机无法学习 XOR 逻辑函数，导致神经网络研究遭受重大挫折。

\begin{figure}[H]
\centering
\includegraphics[width=0.95\textwidth]{slides-images/slide-035.jpg}
\caption{反向传播论文与历史时间线：左为 Rumelhart、Hinton 和 Williams 1986年的开创性论文；右为 Frank Rosenblatt 约1957年与感知机硬件的合影\protect\footnotemark}
\end{figure}
\footnotetext{Slides 第36页（Part 1）。}

\subsubsection{Neocognitron（1980）：第一个``深度''网络}

日本研究者 Fukushima 在 1980 年提出了 \textbf{Neocognitron}——一个约5--7层的神经网络，直接受到 Hubel \& Wiesel 视觉通路研究的启发：

\begin{itemize}
    \item 浅层使用卷积（convolution）函数处理简单特征
    \item 深层汇聚（pooling）浅层信息，处理更复杂的特征
    \item 可以识别手写数字和字母
\end{itemize}

\begin{warningbox}{Neocognitron 的局限}
Neocognitron 的所有参数都是\textbf{手工设计}的——Fukushima 需要一丝不苟地调整数百个参数才能让网络工作。它是一项杰出的工程成就，但缺乏系统化的学习机制，无法扩展到更复杂的任务。
\end{warningbox}

\subsection{反向传播：分水岭时刻（1986）}

\textbf{1986年}，David Rumelhart、Geoffrey Hinton 和 Ronald Williams 发表了具有里程碑意义的论文``Learning Representations by Back-propagating Errors''。反向传播（Backpropagation）基于微积分的链式法则，提供了一种数学上严格的学习规则：

\begin{enumerate}
    \item 将输入通过网络前向传播，得到输出
    \item 计算输出与正确答案之间的误差
    \item 将误差信息\textbf{反向传播}到网络的每一层
    \item 根据梯度信息更新所有参数
\end{enumerate}

\begin{importantbox}{反向传播的革命性意义}
反向传播的出现意味着：
\begin{itemize}
    \item 不再需要\textbf{手工调参}——网络可以通过数据自动学习参数
    \item 理论上可以训练\textbf{任意深度}的网络
    \item 为后续所有深度学习方法奠定了计算基础
\end{itemize}
反向传播至今仍是几乎所有深度学习模型训练的核心算法。
\end{importantbox}

\subsection{LeNet 与早期应用（1998）}

Yann LeCun 在 1990 年代于贝尔实验室创建了 \textbf{LeNet}——大约7层的卷积神经网络，结合反向传播进行端到端训练，成功应用于手写数字和字母识别。LeNet 被部署到美国部分邮局和银行，用于自动读取支票上的数字。

\begin{figure}[H]
\centering
\includegraphics[width=0.95\textwidth]{slides-images/slide-038.jpg}
\caption{LeNet：Yann LeCun 在贝尔实验室开发的卷积神经网络，用于手写字符识别，是最早的商用深度学习系统之一\protect\footnotemark}
\end{figure}
\footnotetext{Slides 第39页（Part 1）。}

然而，LeNet 之后神经网络的发展遭遇了瓶颈：虽然在简单任务（手写数字）上表现优秀，但在复杂的自然图像识别任务（识别猫、狗、椅子、花朵等）上却完全不行。

\subsection{数据的觉醒：ImageNet（2009）}

\begin{figure}[H]
\centering
\includegraphics[width=0.95\textwidth]{slides-images/slide-042.jpg}
\caption{ImageNet 大规模视觉识别挑战赛：包含 1000 个物体类别、超过 143 万张图像\protect\footnotemark}
\end{figure}
\footnotetext{Slides 第43页（Part 1）。}

Fei-Fei Li 和她的团队认识到，整个领域都\textbf{低估了数据的重要性}。神经网络是高容量（high-capacity）模型，需要大量数据才能学会泛化；缺乏数据不仅是不方便，更是一个根本性的数学问题——涉及到泛化和过拟合的深层原理。

\begin{knowledgebox}{ImageNet 的诞生}
\begin{itemize}
    \item 从近10亿张互联网图像中筛选清洗，最终得到 \textbf{1500万}张标注图像
    \item 覆盖 \textbf{22,000} 个物体类别（参考了认知心理学文献，这大约是人类在早期生活中学会识别的类别数量级）
    \item 其子集（100万+图像、1000个类别）被用于举办国际大赛——\textbf{ImageNet Large Scale Visual Recognition Challenge (ILSVRC)}
\end{itemize}
\end{knowledgebox}

\begin{warningbox}{数据是深度学习的一等公民}
在深度学习的早期，大多数研究者只关注网络架构，忽视了数据的重要性。Fei-Fei Li 的洞察是：\textbf{数据和算法一样重要}。高容量模型如果没有足够的数据支撑，必然会过拟合。这一认识对于后来 LLM 时代``数据为王''的理念具有深远的先见性。
\end{warningbox}

\subsection{AlexNet：深度学习革命的引爆点（2012）}

\begin{figure}[H]
\centering
\includegraphics[width=0.95\textwidth]{slides-images/slide-043.jpg}
\caption{ImageNet 挑战赛历年错误率：从 2010 年的 28.2\% 到 2012 年 AlexNet 的 16.4\%（错误率骤降近一半），最终到 2017 年超越人类水平（5.1\%）\protect\footnotemark}
\end{figure}
\footnotetext{Slides 第44页（Part 1）。}

2012 年，Jeff Hinton 和他的学生 Alex Krizhevsky、Ilya Sutskever 参加 ImageNet 挑战赛，使用卷积神经网络（后称 \textbf{AlexNet}），\textbf{将错误率从上一年的 25.8\% 降低到 16.4\%}——几乎减半。

\begin{figure}[H]
\centering
\includegraphics[width=0.95\textwidth]{slides-images/slide-045.jpg}
\caption{AlexNet 的架构：与32年前的 Neocognitron 在结构上并无本质区别，关键差异在于\textbf{反向传播}提供了学习能力，\textbf{大数据}提供了泛化能力\protect\footnotemark}
\end{figure}
\footnotetext{Slides 第46页（Part 1）。}

\begin{importantbox}{从 Neocognitron 到 AlexNet：两个关键突破}
AlexNet 的架构与32年前的 Neocognitron 惊人地相似。它们之间的关键差异在于：
\begin{enumerate}
    \item \textbf{反向传播}（1986）：数学上严格的学习规则，取代了手工调参
    \item \textbf{大规模数据}（ImageNet，2009）：百万级标注图像，让高容量模型能够学会泛化
\end{enumerate}
许多人将 2012 年视为现代 AI 的\textbf{诞生年}或深度学习革命的起点。
\end{importantbox}

\subsection{2012 年后的深度学习爆发}

AlexNet 之后，各种更强大的架构相继诞生，在 ImageNet 挑战赛中持续刷新纪录：

\begin{table}[H]
\centering
\begin{tabular}{lccc}
\toprule
\textbf{年份} & \textbf{模型} & \textbf{Top-5 错误率} & \textbf{关键创新} \\
\midrule
2012 & AlexNet & 16.4\% & CNN + GPU 训练 \\
2013 & ZFNet & 11.7\% & 更好的可视化和调优 \\
2014 & VGGNet / GoogLeNet & 7.3\% / 6.7\% & 更深的网络 / Inception 模块 \\
2015 & ResNet & 3.6\% & 残差连接，152层 \\
2017 & SENet & 2.3\% & 超越人类水平（5.1\%） \\
\bottomrule
\end{tabular}
\caption{ImageNet 挑战赛关键年份的模型和错误率}
\end{table}

\begin{figure}[H]
\centering
\includegraphics[width=0.95\textwidth]{slides-images/slide-047.jpg}
\caption{深度学习时代的里程碑架构：AlexNet 之后出现了 VGGNet、GoogLeNet、ResNet 等一系列影响深远的网络\protect\footnotemark}
\end{figure}
\footnotetext{Slides 第48页（Part 1）。}

\subsection{本章小结}

深度学习的崛起是一个长达数十年的过程：从感知机（1958）到反向传播（1986）到 LeNet（1998），再到 ImageNet（2009）和 AlexNet（2012）。2012 年的 ImageNet 挑战赛是一个历史性的转折点，两大要素的汇聚——反向传播算法和大规模标注数据——彻底释放了神经网络的潜力，开启了深度学习革命。

%% ========== 深度学习的广泛影响 ==========

